\documentclass[11pt]{article}

\usepackage[a4paper,margin=2cm]{geometry}
\usepackage{graphicx}
\usepackage{amsmath}
\usepackage{xcolor}
\usepackage{enumitem}
\usepackage{booktabs}
\usepackage{titlesec}
\usepackage{longtable}

\graphicspath{{../figures/}}

\definecolor{accent}{RGB}{0,90,160}
\definecolor{ink}{RGB}{28,28,28}
\definecolor{theorybg}{RGB}{234,241,248}
\color{ink}

\setlist[itemize]{leftmargin=1.4em,itemsep=2pt,topsep=2pt,parsep=0pt}
\titleformat{\section}{\large\bfseries\color{accent}}{}{0pt}{}
\titleformat{\subsection}{\normalsize\bfseries\color{accent}}{}{0pt}{}
\titlespacing{\section}{0pt}{4pt}{3pt}
\titlespacing{\subsection}{0pt}{5pt}{1pt}
\pagestyle{plain}
\setlength{\parindent}{0pt}

% ---- helpers ------------------------------------------------------------
\newcommand{\figblock}[1]{%
  \begin{center}\includegraphics[width=0.84\linewidth]{#1}\end{center}}
\newcommand{\theorybox}[1]{%
  \vspace{1pt}\noindent
  \colorbox{theorybg}{\begin{minipage}{\dimexpr\linewidth-2\fboxsep\relax}
  \small\textbf{\textcolor{accent}{Theory in brief.}}\ #1
  \end{minipage}}\vspace{3pt}}
\newcommand{\shows}{\subsection*{\textcolor{accent}{1.\ \ What the plot shows}}\small}
\newcommand{\method}{\subsection*{\textcolor{accent}{2.\ \ Method --- how the data was obtained}}\small}
\newcommand{\analysis}{\subsection*{\textcolor{accent}{3.\ \ Analysis, deviations and conclusion}}\small}

\begin{document}

% =========================================================================
% TITLE PAGE
% =========================================================================
\begin{titlepage}
\centering
\vspace*{2.0cm}
{\Huge\bfseries\color{accent} Robustness Analysis of the\\[6pt]
SSH Topolectrical Circuit\par}
\vspace{10pt}
{\Large An independent, from-scratch reconstruction\par}
\vspace{1.2cm}
{\large Pristine baseline \;$\cdot$\; Single impurities \;$\cdot$\;
Domain walls \;$\cdot$\; Monte~Carlo disorder\par}
\vspace{1.8cm}

\begin{minipage}{0.86\linewidth}
\color{ink}\large
This report rebuilds the SSH topolectrical-circuit robustness study
\emph{from the ground up}. Every circuit was re-simulated and every number
re-derived with a purpose-written nodal-admittance solver; no data from any
previous run was reused. The circuit uses the \textbf{updated component
values}:
\vspace{6pt}
\begin{center}
\begin{tabular}{lll}
\toprule
Intracell capacitor $C_1$ & $2.2\ \mathrm{nF}$ & series resistance $1\ \Omega$\\
Intercell capacitor $C_2$ & $3.3\ \mathrm{nF}$ & series resistance $1\ \Omega$\\
Grounded inductor $L$     & $2200\ \mu\mathrm{H}$ & series resistance $1\ \Omega$\\
\bottomrule
\end{tabular}
\end{center}
\vspace{6pt}
Every result is checked against analytic theory as it is produced, and every
deviation from the ideal expectation is explained. The headline finding: the
topological boundary resonance at \textbf{45.74~kHz} survives realistic
component disorder up to $\pm10\%$ essentially intact.
\end{minipage}

\vfill
{\large 20 May 2026\par}
\end{titlepage}

% =========================================================================
% METHOD PAGE
% =========================================================================
\section*{The circuit, the method, and how it was checked}

\textbf{The circuit.}
The Su--Schrieffer--Heeger (SSH) topolectrical circuit is a chain of
$N=8$ unit cells, giving \textbf{16 nodes}. Every node is tied to ground
through an inductor $L$. Neighbouring nodes are linked by capacitors that
\emph{alternate} between a small intracell value $C_1$ and a larger
intercell value $C_2$; a grounded boundary capacitor $C_b=C_2$ closes the
left edge, and an AC voltage source drives node~1. Because $C_1<C_2$ the
circuit is in the \emph{topologically non-trivial} phase, which hosts a
protected edge state --- the electrical analogue of the SSH model from
condensed-matter physics. Technical terms are defined in the
glossary (Appendix~C).

\medskip
\textbf{The simulation engine.}
A from-scratch solver (\texttt{ssh\_circuit.py}) assembles the
$16\times16$ complex nodal-admittance matrix of the circuit at each
frequency and solves the linear system for every node voltage --- this is
\emph{exactly} the calculation an industry SPICE \texttt{.ac} analysis
performs, so the output is SPICE-accurate by construction while remaining
transparent and reproducible. The assembly and solve are vectorised over
frequency, so the Monte~Carlo of $>50{,}000$ circuits runs in about a
minute. LTSpice netlists of the key circuits are also exported (folder
\texttt{netlists/}) for bench reproduction.

\medskip
\textbf{The checking philosophy.}
Every figure is cross-checked against an independent theory prediction
\emph{as it is produced}, and every gap between simulation and the
idealised expectation is traced to its cause. The core checks all passed:

\begin{center}\small
\begin{tabular}{lll}
\toprule
\textbf{Quantity} & \textbf{Analytic theory} & \textbf{This simulation}\\
\midrule
TBR frequency $f_R=1/(2\pi\sqrt{L(C_1{+}C_2)})$ & $45.754$ kHz
   & $45.742$ kHz \ (0.025\% --- finite-size)\\
Edge-state decay length $\xi=1/\ln(C_2/C_1)$ & $2.4663$ cells
   & $2.4656$ cells \ (0.03\%)\\
Per-cell voltage ratio $=C_1/C_2$ & $0.66667$ & $0.66658$\\
Resonance from Laplacian eigenvalue & $45.754$ kHz & $45.742$ kHz\\
Impurity sign ($\delta>0 \Rightarrow f$ falls) & predicted & confirmed\\
Monte-Carlo spread linear in $\sigma$ & ratio $=1.00$
   & ratio $=1.01,\ 1.03$\\
\bottomrule
\end{tabular}
\end{center}

\medskip
\textbf{How to read this report.}
The next 16 pages each carry \emph{one figure}. Below it sits a shaded
\textbf{Theory in brief} box, then three numbered blocks: \textbf{(1)} what
the plot shows, \textbf{(2)} how the data was obtained, and \textbf{(3)} a
quantitative analysis that states the conclusion and explains every
deviation from ideal behaviour. The figures move through four parts: the
pristine baseline (Figs.~1--3), single impurities (Figs.~4--5), domain
walls (Figs.~6--8), and the Monte~Carlo disorder study (Figs.~9--16).
Appendix~A lists the sources, Appendix~B tabulates every numerical result
against theory, and Appendix~C is a glossary of technical terms.

% =========================================================================
% FIGURE 1
% =========================================================================
\clearpage
\section{Figure 1\quad Pristine circuit: impedance spectrum and the TBR}
\figblock{fig01_pristine_impedance.png}
\theorybox{Any inductor--capacitor loop ``rings'' at a resonant frequency
$f=1/(2\pi\sqrt{LC})$. The 16-node SSH chain supports 16 such modes, whose
frequencies cluster into two \emph{bulk bands} separated by an empty
\emph{band gap}. Topology guarantees one extra mode \emph{inside} that
gap --- the topological boundary resonance (TBR) --- bound to the edge of
the chain, at $f_R=1/(2\pi\sqrt{L(C_1+C_2)})$.}

\shows
\begin{itemize}
\item The impedance $|Z_{\rm in}|$ that the AC source ``feels'' at node~1,
      swept from 20 to 120~kHz (linear frequency axis, logarithmic
      impedance axis).
\item One tall, sharp peak at \textbf{45.74~kHz} --- the TBR --- standing
      alone in the white band gap; the smaller peaks inside the grey strips
      are ordinary bulk resonances.
\item The orange dashed line is the textbook prediction $f_R=45.75$~kHz.
\end{itemize}

\method
\begin{itemize}
\item The 8-cell, 16-node circuit was built from scratch with the updated
      parts ($C_1=2.2$~nF, $C_2=3.3$~nF, $L=2200\ \mu$H, every component
      carrying $1\ \Omega$ series resistance).
\item At each of 1600 frequencies the $16\times16$ nodal-admittance matrix
      was assembled and solved for the node voltages; $|Z_{\rm in}|$ is the
      source voltage divided by the source current.
\item The peak was located with a coarse-then-fine sweep and a parabolic
      refinement.
\end{itemize}

\analysis
\begin{itemize}
\item \textbf{Quantitative:} the simulated peak is $45.742$~kHz; the
      textbook formula gives $45.754$~kHz; the independent eigenvalue
      method (Fig.~3) gives $45.742$~kHz. Peak quality factor $Q=892$,
      width only $51$~Hz. The gap runs $41.8$--$51.2$~kHz (9.4~kHz wide)
      and the TBR sits 42\% of the way up it.
\item \textbf{Deviation explained:} the 0.025\% gap between the simulated
      peak and the textbook formula is \emph{not} an error. The formula
      assumes an infinitely long chain; in a finite 8-cell chain the
      left-edge and right-edge modes overlap by $(C_1/C_2)^8=3.9\%$, and
      this slight hybridisation pulls the mode down by $\approx0.025\%$.
      The eigenvalue method, which also uses the finite chain, agrees with
      the simulation to 5~significant figures --- confirming the cause.
\item \textbf{Qualitative:} one dominant peak in the gap, bulk peaks an
      order of magnitude smaller; the grey/white band pattern matches SSH
      band theory.
\item \textbf{Conclusion:} the circuit and solver are correct, and the TBR
      exists as a sharp, isolated, easily measurable resonance --- the
      foundation for everything that follows.
\end{itemize}

% =========================================================================
% FIGURE 2
% =========================================================================
\clearpage
\section{Figure 2\quad Pristine circuit: edge-state voltage profile}
\figblock{fig02_pristine_voltage.png}
\theorybox{An edge state is a resonance whose voltage is concentrated at one
end of the chain and dies away \emph{geometrically} into the bulk --- by a
fixed factor $C_1/C_2$ per unit cell. Its decay length is
$\xi=1/\ln(C_2/C_1)$. Chiral symmetry forces the mode to live almost
entirely on one of the two sublattices (here the odd-numbered nodes).}

\shows
\begin{itemize}
\item The voltage at each of the 16 nodes (divided by the source voltage),
      driven at the TBR, on a logarithmic axis.
\item \textbf{Odd nodes} (blue): 1.0 at the driven edge, falling
      exponentially rightwards. \textbf{Even nodes} (green): near-zero at
      the left, rising towards the right edge --- the chain's second edge
      state.
\item The orange dashed line is an exponential fit to the odd nodes.
\end{itemize}

\method
\begin{itemize}
\item The circuit was solved once, exactly at the TBR frequency, and the
      magnitude of all 16 node voltages recorded.
\item A straight line was fitted to the logarithm of the odd-node
      voltages; its slope gives the decay length $\xi$.
\end{itemize}

\analysis
\begin{itemize}
\item \textbf{Quantitative:} the fitted decay length is $\xi=2.4656$~cells
      against the theory $\xi=1/\ln(C_2/C_1)=2.4663$ --- a 0.03\% match.
      The per-cell voltage ratio is $0.66658$ versus $C_1/C_2=0.66667$.
      The exponential fit has $R^2=0.999\,999\,98$.
\item \textbf{Qualitative:} the profile is a clean two-edge structure ---
      odd nodes decay from the left, even nodes rise to the right ---
      exactly the chiral fingerprint of the SSH model.
\item \textbf{Deviations:} none beyond the fifth decimal place; the tiny
      residual is the same finite-size hybridisation seen in Fig.~1.
\item \textbf{Conclusion:} the TBR is a genuine, exponentially localised
      topological edge state. This profile is the reference shape against
      which every later disorder result is measured.
\end{itemize}

% =========================================================================
% FIGURE 3
% =========================================================================
\clearpage
\section{Figure 3\quad Pristine resonance spectrum: the TBR isolated in the gap}
\figblock{fig03_resonance_spectrum.png}
\theorybox{Each resonance of the circuit corresponds to one eigenvalue
$\mu$ of its lossless capacitance Laplacian, through
$f=1/(2\pi\sqrt{L\mu})$. The bulk eigenvalues form two bands. The
\emph{bulk--boundary correspondence} theorem guarantees that whenever the
chain is in the topological phase ($C_1<C_2$) exactly \emph{one} eigenvalue
falls inside the gap --- the TBR.}

\shows
\begin{itemize}
\item All 16 resonance frequencies of the circuit, sorted low-to-high
      against a mode index.
\item Grey bands mark the two bulk bands predicted by SSH theory; the white
      strip is the band gap.
\item The blue star is the TBR --- alone in the gap, on the analytic
      $f_R$ line.
\end{itemize}

\method
\begin{itemize}
\item The eigenvalues $\mu$ of the lossless capacitance Laplacian $W$ were
      computed; each gives a resonance $f=1/(2\pi\sqrt{L\mu})$.
\item The bulk-band edges were taken from the tight-binding SSH model.
\end{itemize}

\analysis
\begin{itemize}
\item \textbf{Quantitative:} the 16 modes split as 7 in the lower band,
      \textbf{1 in the gap}, and 8 in the upper band. The gap mode lies at
      $45.742$~kHz, matching the impedance peak of Fig.~1 to 5~figures.
\item \textbf{Qualitative:} this is a third, fully independent route to
      $f_R$ (after the impedance sweep and the textbook formula) and it
      agrees with both --- a triple cross-check that the model is sound.
\item \textbf{Deviations:} the gap mode sits $0.025\%$ below the formula's
      $45.754$~kHz, again the finite-size hybridisation; it is not exactly
      mid-gap (42\% up) because the two bulk bands are not symmetric about
      the gap centre.
\item \textbf{Conclusion:} there is exactly one protected mode in the gap,
      precisely as the bulk--boundary theorem requires. The width of the
      gap is the ``armour'' that the disorder studies later test.
\end{itemize}

% =========================================================================
% FIGURE 4
% =========================================================================
\clearpage
\section{Figure 4\quad Single impurity: TBR displacement vs.\ position and strength}
\figblock{fig04_single_impurity_shift.png}
\theorybox{Changing a capacitor changes the resonant frequency, because $f$
depends on the capacitances. To \emph{first order} the shift is
proportional to the defect strength $\delta$ and to how much of the edge
mode sits in that cell. But $f\propto1/\sqrt{\text{capacitance}}$ is a
\emph{curved} function, so equal $\pm\delta$ changes give slightly unequal
shifts --- the response is not perfectly linear.}

\shows
\begin{itemize}
\item How the TBR frequency (ratio $f_{\rm peak}/f_R$) moves when one unit
      cell is made impure --- both its capacitors scaled by $(1+\delta)$.
\item Solid lines: simulation for $\delta=\pm0.1,\pm0.3$ across cells
      2--7. Dashed lines: first-order perturbation theory.
\item Positive $\delta$ lowers the resonance, negative $\delta$ raises it;
      the effect is largest at cell~2 and negligible by cell~7.
\end{itemize}

\method
\begin{itemize}
\item For every (cell, $\delta$) pair the two capacitors of that cell were
      rescaled, the impedance sweep re-run and the new peak extracted ---
      24 independent simulations.
\item The theory line is the first-order change of the Laplacian
      eigenvalue computed from the undisturbed edge-mode shape.
\end{itemize}

\analysis
\begin{itemize}
\item \textbf{Quantitative:} the largest shifts are at cell~2 ---
      $+2.57\%$ for $\delta=-0.3$ and $-2.40\%$ for $\delta=+0.3$. The
      ``leverage'' of cell~2 over cell~7 is a factor of 51, tracking the
      edge mode's own amplitude.
\item \textbf{Asymmetry explained:} the $+2.57\%$ vs.\ $2.40\%$ inequality
      (a factor 1.07) is the curvature of $f\propto1/\sqrt{\;}$: shrinking
      a capacitor raises $f$ slightly more than enlarging it lowers $f$.
\item \textbf{Deviation from theory:} simulation and first-order theory
      agree to better than $0.1\%$ everywhere; the only visible gap is at
      cell~2, $\delta=\pm0.3$, the expected onset of non-linearity for the
      strongest single defect.
\item \textbf{Conclusion:} a local defect shifts the TBR predictably and
      by at most 2.6\% --- it perturbs the edge state but cannot destroy
      it. The position dependence is a direct map of where the edge mode
      lives.
\end{itemize}

% =========================================================================
% FIGURE 5
% =========================================================================
\clearpage
\section{Figure 5\quad Single impurity in cell 3: voltage profile}
\figblock{fig05_single_impurity_profile.png}
\theorybox{The driven edge mode is pinned at node~1 and decays rightwards,
so a defect \emph{downstream} cannot change what lies \emph{upstream} of
it. A single-cell defect also weakly breaks chiral symmetry locally, which
lets a little voltage leak onto the B sublattice (even nodes) that the
pristine edge state leaves almost empty.}

\shows
\begin{itemize}
\item The voltage profile with cell~3 impure ($\delta=\pm0.3$) versus
      pristine. The two sublattices are drawn as \emph{separate} curves
      (solid = A/odd, dashed = B/even) --- no zig-zag.
\item The A-sublattice curves lie on top of each other up to the impurity,
      then separate; the B-sublattice is lifted near the defect.
\end{itemize}

\method
\begin{itemize}
\item As Figure~2, but with cell-3 capacitors rescaled; each profile is
      read at \emph{that} circuit's own (shifted) resonance, so the
      comparison is like-for-like.
\item The deviation from the pristine profile was measured node by node.
\end{itemize}

\analysis
\begin{itemize}
\item \textbf{Quantitative:} on the dominant A sublattice the upstream
      nodes (N1, N3) change by $\le0.32\%$ --- effectively frozen ---
      while downstream the change reaches $6.0\%$ for $\delta=+0.3$ and
      $28.6\%$ for $\delta=-0.3$.
\item \textbf{Asymmetry explained:} $\delta=-0.3$ reshapes the chain far
      more than $\delta=+0.3$ because it \emph{weakens} the coupling out
      of cell~3, partly decoupling the downstream chain; $\delta=+0.3$
      strengthens it.
\item \textbf{The B sublattice:} pristine even-node voltages are tiny
      ($\le0.039$); the defect lifts them up to $0.12$ near itself --- a
      real chiral-symmetry-breaking leak, but still far below the
      A-sublattice values, so the edge state stays A-dominated.
\item \textbf{Conclusion:} the impurity reshapes the chain only
      \emph{from itself onwards}; the protected edge mode upstream is
      untouched and the state remains edge-localized.
\end{itemize}

% =========================================================================
% FIGURE 6
% =========================================================================
\clearpage
\section{Figure 6\quad AB impurity pair: TBR displacement vs.\ $\gamma$}
\figblock{fig06_AB_pair_sweep.png}
\theorybox{Two adjacent defects --- an ``AB pair'' --- change three
capacitors at once. The resonance shift again follows
$f\propto1/\sqrt{\mu}$ with the eigenvalue $\mu$ \emph{linear} in the
defect strengths. A linear input through a square-root gives a
\emph{curved} output: first-order theory is only the straight tangent
line, exact near $\gamma=0$ and drifting away for strong defects.}

\shows
\begin{itemize}
\item The TBR shift when an AB pair sits at cell~3, sweeping the second
      site's strength $\gamma$ for three fixed values of the first,
      $\delta$.
\item Solid = simulation; dashed = first-order perturbation theory.
      $f_{\rm peak}/f_R$ falls as $\gamma$ rises, along gently
      \emph{curved} lines.
\end{itemize}

\method
\begin{itemize}
\item An AB pair changes three capacitors (two intercell, one intracell
      set to the average of $\delta$ and $\gamma$), per Anzurez \textit{et
      al.}\ (2025).
\item $\gamma$ was swept over 25 values; the simulation was also compared
      to the \emph{exact} lossless-eigenvalue prediction.
\end{itemize}

\analysis
\begin{itemize}
\item \textbf{Quantitative:} the shift spans $3.8\%$. The
      $\delta=0$ curve has a measured curvature of $+0.018$, and $-\gamma$
      shifts the resonance $1.40\times$ more than $+\gamma$ --- both direct
      consequences of the $1/\sqrt{\;}$ convexity.
\item \textbf{Deviation from theory explained:} the simulation departs
      from the dashed first-order line by up to $0.91\%$. This is
      \emph{not} a simulation error: the same circuits reproduce the exact
      lossless-eigenvalue prediction to $0.0004\%$. The gap is entirely the
      error of \emph{linearised} theory applied to a strong, three-capacitor
      perturbation.
\item \textbf{Asymmetry:} the three $\delta$ curves are unequally
      spaced --- the $\delta=-0.2$ curve sits further from $\delta=0$ than
      $\delta=+0.2$ does, the same convexity acting on $\delta$.
\item \textbf{Conclusion:} even a clustered two-site defect moves the TBR
      by only a few percent; the behaviour is fully understood as a curved
      (not linear) but smooth and monotonic response.
\end{itemize}

% =========================================================================
% FIGURE 7
% =========================================================================
\clearpage
\section{Figure 7\quad BA impurity pair: TBR displacement vs.\ $\gamma$}
\figblock{fig07_BA_pair_sweep.png}
\theorybox{A BA pair is an AB pair with the two impurity sites in the
opposite order: it changes two intracell capacitors and one intercell
capacitor (the mirror of the AB case). The same $f\propto1/\sqrt{\mu}$
physics applies, so the response is again curved rather than straight.}

\shows
\begin{itemize}
\item The same experiment as Figure~6 but for a BA pair.
\item The resonance again moves smoothly with $\gamma$, along
      \emph{visibly curved} lines, over a noticeably \emph{narrower} range
      than the AB pair.
\end{itemize}

\method
\begin{itemize}
\item A BA pair changes three capacitors in the mirror pattern of the AB
      pair; $\gamma$ swept as before, with first-order theory dashed and an
      exact-eigenvalue cross-check.
\end{itemize}

\analysis
\begin{itemize}
\item \textbf{The curvature explained:} every line bows because
      $f_{\rm peak}/f_R = 1/\sqrt{1+(\text{linear in }\gamma)}$ --- a
      square root, not a straight line. The fitted curvature is $+0.019$,
      slightly stronger than the AB pair's.
\item \textbf{Quantitative:} the BA shift spans only $2.5\%$, against the
      AB pair's $3.8\%$ --- BA disturbs the TBR less. The $-\gamma$/$+\gamma$
      asymmetry is $1.69\times$.
\item \textbf{Deviation from theory:} the simulation hugs first-order
      theory far better than the AB case (max gap $0.27\%$ vs.\ $0.91\%$)
      and reproduces the exact eigenvalue to $0.0002\%$ --- the curvature
      is genuine physics, captured exactly by the solver.
\item \textbf{Conclusion:} defect \emph{ordering} matters --- BA perturbs
      the TBR roughly $1.5\times$ less than AB --- and the curved response
      is the expected square-root law, not noise.
\end{itemize}

% =========================================================================
% FIGURE 8
% =========================================================================
\clearpage
\section{Figure 8\quad Domain wall: voltage profiles for AB and BA pairs}
\figblock{fig08_domain_wall_profile.png}
\theorybox{An impurity pair locally reverses the chain's strong/weak bond
pattern --- a \emph{domain wall}. In the SSH model a domain wall binds a
localised mid-gap state (a ``soliton''). In the circuit this appears as
voltage piling up on the normally-empty B sublattice near the pair.}

\shows
\begin{itemize}
\item Voltage profiles for a strong AB pair and a strong BA pair
      ($\delta=\gamma=0.3$) versus pristine, sublattices drawn separately.
\item The B sublattice (dashed) is lifted dramatically near the defect ---
      voltage accumulates where the pristine chain had almost none.
\end{itemize}

\method
\begin{itemize}
\item The AB and BA circuits were built and solved at their own
      resonances; the evenness of each profile was scored with an
      inverse-participation-ratio (IPR) over the even nodes.
\end{itemize}

\analysis
\begin{itemize}
\item \textbf{Quantitative:} the pristine B sublattice peaks at only
      $0.039$; the AB and BA pairs lift it to $0.16$ --- a fourfold
      pile-up of charge around the defect, the circuit signature of the
      bound domain-wall state.
\item \textbf{AB vs.\ BA:} the even-node IPR is $0.57$ for AB but only
      $0.38$ for BA --- BA spreads the domain-wall charge far more evenly,
      confirming the published expectation.
\item \textbf{Qualitative:} the A sublattice keeps its overall exponential
      shape; the new physics is entirely on the B sublattice, exactly as
      a chiral-symmetry-breaking domain wall predicts.
\item \textbf{Conclusion:} an impurity pair behaves as a controllable
      domain wall that binds extra charge --- the circuit analogue of the
      SSH soliton --- and its shape can be tuned by the defect ordering.
\end{itemize}

% =========================================================================
% FIGURE 9
% =========================================================================
\clearpage
\section{Figure 9\quad Monte Carlo: TBR frequency distribution vs.\ disorder}
\figblock{fig09_mc_fpeak_hist.png}
\theorybox{Real capacitors differ from their nominal value by a
manufacturing tolerance. Modelling each capacitor as
nominal$\,\times(1+\varepsilon)$ with $\varepsilon$ a random Gaussian of
width $\sigma$, and repeating thousands of times, turns the single TBR into
a \emph{distribution}. A topologically protected mode should broaden but
never vanish.}

\shows
\begin{itemize}
\item The spread of the TBR frequency across \textbf{8000 random circuits}
      at each of three disorder levels, $\sigma=1\%,5\%,10\%$.
\item At $1\%$ a tight spike; at $10\%$ a much broader but still
      \emph{single} peak, centred on the pristine value.
\end{itemize}

\method
\begin{itemize}
\item In each circuit every capacitor was independently jittered by a
      Gaussian of width $\sigma$; the TBR of each of the 24{,}000 circuits
      was extracted and binned.
\item Each realisation uses a fixed random seed, so the run is exactly
      reproducible.
\end{itemize}

\analysis
\begin{itemize}
\item \textbf{Quantitative:} the distribution full-width grows
      $0.24\to1.20\to2.45$~kHz as $\sigma$ goes $1\to5\to10\%$ --- roughly
      proportional to $\sigma$ (analysed in Fig.~10).
\item \textbf{Qualitative:} the peak stays single and symmetric at every
      level; it never collapses, splits, or develops a second peak ---
      all of which would signal destruction of the edge mode.
\item \textbf{Deviations:} none qualitative; $10\%$ disorder is a
      quantitative broadening, not a breakdown.
\item \textbf{Conclusion:} the TBR persists as a clear, identifiable
      resonance at every realistic component tolerance --- the first
      direct evidence of topological robustness.
\end{itemize}

% =========================================================================
% FIGURE 10
% =========================================================================
\clearpage
\section{Figure 10\quad Monte Carlo: TBR frequency --- mean and spread vs.\ $\sigma$}
\figblock{fig10_mc_fpeak_vs_sigma.png}
\theorybox{To first order a random capacitor change shifts $f$ up or down
with equal chance, so the ensemble mean should stay at $f_R$ and the spread
should grow $\propto\sigma$. A small \emph{upward} bias is nonetheless
expected: $f\propto1/\sqrt{\mu}$ is a convex function, and the average of a
convex function exceeds the function of the average (Jensen's inequality).}

\shows
\begin{itemize}
\item A condensed view of the Figure-9 data: the ensemble mean TBR
      frequency (point) with an error bar of $\pm1$ standard deviation.
\item The red dashed line is the pristine resonance.
\end{itemize}

\method
\begin{itemize}
\item For each $\sigma$ the mean and standard deviation were taken over the
      8000 measured frequencies; the error bar is that standard deviation,
      the genuine physical spread of the ensemble.
\end{itemize}

\analysis
\begin{itemize}
\item \textbf{Quantitative:} the spread grows almost perfectly linearly ---
      $0.10\to0.51\to1.04$~kHz, i.e.\ $1.01\times$ and $1.03\times$ the
      linear expectation at $5\%$ and $10\%$.
\item \textbf{The bias explained:} the mean is not exactly on the pristine
      line --- it drifts up by $+0.003\%$, $+0.09\%$, $+0.34\%$ at
      $\sigma=1,5,10\%$. This is the predicted Jensen / square-root
      convexity effect: it grows like $\sigma^2$ (note $0.34/0.09\approx4$
      for a doubling of $\sigma$), confirming the explanation.
\item \textbf{Conclusion:} disorder adds \emph{scatter}, and only a tiny,
      fully-understood second-order \emph{bias}; the resonance is not
      systematically destabilised. Even the worst single circuit out of
      8000 stayed within $-15\%/+12\%$ of $f_R$.
\end{itemize}

% =========================================================================
% FIGURE 11
% =========================================================================
\clearpage
\section{Figure 11\quad Monte Carlo: peak impedance vs.\ disorder}
\figblock{fig11_mc_zpeak_vs_sigma.png}
\theorybox{The peak impedance is largest when every component sits at its
matched nominal value. \emph{Any} disorder --- a capacitor too big
\emph{or} too small --- detunes that match, and detuning can only lower the
peak. So the mean peak height should drift \emph{down}, by an amount that
grows quadratically ($\propto\sigma^2$) with the disorder.}

\shows
\begin{itemize}
\item The height of the resonance peak, $|Z|_{\rm peak}$, as an ensemble
      mean with a $\pm1$ s.d.\ error bar at each $\sigma$.
\item The mean stays nearly flat; the spread grows with $\sigma$.
\end{itemize}

\method
\begin{itemize}
\item For each of the 24{,}000 disordered circuits the peak impedance was
      recorded alongside the peak frequency; the ensemble mean and standard
      deviation are plotted.
\end{itemize}

\analysis
\begin{itemize}
\item \textbf{Quantitative:} the mean peak height is
      $146.3\to146.2\to145.0\to141.0$~k$\Omega$ for
      $\sigma=0,1,5,10\%$ --- a fall of only $3.6\%$ at the worst case.
\item \textbf{The bias explained:} the mean drifts \emph{down} (unlike the
      frequency, which drifts up), by $-0.04\%,-0.91\%,-3.63\%$. The
      $\sigma^2$ scaling is clear --- a $2\times$ rise in $\sigma$
      from $5\%$ to $10\%$ gives a $\approx4\times$ larger drop --- exactly
      the detuning argument above.
\item \textbf{Qualitative:} a tall peak means a strong edge mode still
      coupled cleanly to the probe; if disorder were hybridising the edge
      state into the bulk, the height would collapse, not merely dip 4\%.
\item \textbf{Conclusion:} the resonance keeps its \emph{strength}, not
      just its position; the small, $\sigma^2$ loss of height is the
      expected detuning effect and nothing more.
\end{itemize}

% =========================================================================
% FIGURE 12
% =========================================================================
\clearpage
\section{Figure 12\quad Monte Carlo: edge-state localization (IPR) vs.\ disorder}
\figblock{fig12_mc_ipr_hist.png}
\theorybox{The inverse participation ratio (IPR) is one number that
measures how concentrated a voltage profile is: high means localized, low
means spread out (the fully-delocalized value here is $1/15=0.067$).
Capacitor disorder perturbs only the \emph{couplings} of the chain and
keeps it bipartite, so the chiral-protected edge mode cannot delocalize.}

\shows
\begin{itemize}
\item The distribution of the IPR over the 8000 circuits at each $\sigma$.
\item Pristine IPR $=0.385$ (black dashed); the fully-delocalized
      reference $1/15=0.067$ (pink dotted) sits far to the left.
\end{itemize}

\method
\begin{itemize}
\item For every disordered circuit the voltage profile was measured and its
      IPR computed over nodes 2--16 (the driven node is excluded so the
      metric is not trivially dominated by it).
\end{itemize}

\analysis
\begin{itemize}
\item \textbf{Quantitative:} the mean IPR falls only slightly with
      disorder --- $0.385\to0.384\to0.376\to0.355$ for
      $\sigma=0,1,5,10\%$ --- staying roughly \emph{five times} the
      delocalized value even at $10\%$.
\item \textbf{Qualitative:} the distributions broaden but remain clustered
      well to the right; not one of the 8000 worst-case circuits comes
      anywhere near the delocalized line.
\item \textbf{Deviation:} the modest downward drift of the mean is the
      same mild ``bulk admixture'' that lengthens $\xi$ in Fig.~13 --- a
      small reshaping, not delocalization.
\item \textbf{Conclusion:} this is the direct numerical proof of
      topological protection --- capacitor disorder cannot delocalize the
      chiral-protected edge mode.
\end{itemize}

% =========================================================================
% FIGURE 13
% =========================================================================
\clearpage
\section{Figure 13\quad Monte Carlo: edge-state decay length vs.\ disorder}
\figblock{fig13_mc_xi_vs_sigma.png}
\theorybox{The decay length $\xi$ is the edge state expressed as a physical
size --- how many cells the voltage takes to fall by a factor $e$.
Disorder mixes a little delocalized \emph{bulk} character into the mode,
which lengthens the average decay, so a small \emph{positive} bias in
$\langle\xi\rangle$ is expected. True delocalization would instead send
$\xi\to\infty$.}

\shows
\begin{itemize}
\item The decay length $\xi$ as an ensemble mean with $\pm1$ s.d.\ error
      bars.
\item The orange dashed line is the theory value $\xi=2.466$ cells.
\end{itemize}

\method
\begin{itemize}
\item For each disordered circuit a straight line was fitted to the
      logarithm of the odd-node voltages; its slope gives $\xi$. The
      ensemble mean and standard deviation are plotted.
\end{itemize}

\analysis
\begin{itemize}
\item \textbf{Quantitative:} the mean $\xi$ rises only gently ---
      $2.466\to2.468\to2.512\to2.686$ cells for $\sigma=0,1,5,10\%$,
      a $+8.9\%$ change at the worst case.
\item \textbf{The bias explained:} the upward drift is the expected
      bulk-admixture effect, and it grows faster than linearly ($+1.9\%$ at
      $5\%$, $+8.9\%$ at $10\%$) because disorder mixes in bulk character
      quadratically. The growing error bars reflect the harder
      log-linear fit on noisier profiles.
\item \textbf{Deviation:} $\xi$ never trends towards infinity --- the
      unmistakable signature that the mode stays localized.
\item \textbf{Conclusion:} the edge state keeps its characteristic size
      to within $\sim9\%$ even at $10\%$ disorder; together with Fig.~12 it
      shows the mode keeps both its shape and its length scale.
\end{itemize}

% =========================================================================
% FIGURE 14
% =========================================================================
\clearpage
\section{Figure 14\quad Monte Carlo: edge-state voltage envelope}
\figblock{fig14_mc_voltage_envelope.png}
\theorybox{If the edge state is genuinely protected, every disordered copy
of it should still be the same exponentially-decaying shape, give or take a
spread. Plotting the ensemble \emph{median} profile together with a
\emph{percentile band} shows that spread directly.}

\shows
\begin{itemize}
\item The A-sublattice (odd-node) voltage profile under disorder: the
      pristine profile (black), the ensemble median at each $\sigma$
      (coloured lines) and the 5\%--95\% spread band (shaded).
\item The medians sit essentially on the pristine line; the bands widen
      with $\sigma$.
\end{itemize}

\method
\begin{itemize}
\item The 8000 measured profiles at each $\sigma$ were stacked; the median
      and the 5th/95th percentiles were taken node by node.
\item Only the dominant A sublattice is shown, so the exponential decay is
      seen cleanly without the odd/even zig-zag.
\end{itemize}

\analysis
\begin{itemize}
\item \textbf{Quantitative:} the three median curves are
      indistinguishable from the pristine exponential. At a mid-chain node
      the $10\%$ band reaches only $\pm50\%$ in voltage --- under a factor
      of two even in the worst case.
\item \textbf{Qualitative:} the characteristic straight-line (exponential)
      decay survives intact; disorder blurs the profile but never erases
      its shape.
\item \textbf{Deviations:} the band is slightly wider towards the far
      (right) end of the chain, where the small voltages are more
      sensitive --- expected, and still bounded.
\item \textbf{Conclusion:} the \emph{whole} edge-state profile, not just
      its summary numbers, is preserved under realistic disorder.
\end{itemize}

% =========================================================================
% FIGURE 15
% =========================================================================
\clearpage
\section{Figure 15\quad Chiral-symmetry test: capacitor vs.\ inductor disorder}
\figblock{fig15_mc_chiral_control.png}
\theorybox{The SSH edge state is protected by \emph{chiral symmetry}.
Capacitor disorder perturbs the off-diagonal \emph{couplings} of the
circuit Laplacian and largely respects that symmetry; inductor disorder
perturbs the on-site (diagonal) terms and \emph{breaks} it. So inductor
disorder should destabilise the resonance more.}

\shows
\begin{itemize}
\item How much the resonance \emph{spreads} (standard deviation of
      $f_{\rm peak}$, as a percentage of $f_R$) for disorder on the
      \emph{capacitors} (blue) versus the \emph{inductors} (pink).
\item Both grow linearly with $\sigma$; the inductor line is clearly
      steeper.
\end{itemize}

\method
\begin{itemize}
\item The Monte~Carlo was run twice --- once jittering only capacitors,
      once only inductors --- 48{,}000 circuits in total.
\item Error bars are bootstrap estimates of the uncertainty on each
      spread.
\end{itemize}

\analysis
\begin{itemize}
\item \textbf{Quantitative:} at $\sigma=10\%$ the frequency spread is
      $2.27\%$ for capacitor disorder but $3.13\%$ for inductor disorder
      --- inductor disorder is $1.38\times$ more destabilising. The mean
      IPR at $10\%$ is also lower for inductor disorder ($0.338$ vs.\
      $0.355$).
\item \textbf{Why the contrast is moderate, not extreme:} capacitor
      disorder is not \emph{perfectly} chiral-symmetric --- it also
      slightly disturbs each node's total grounding capacitance (a
      diagonal term). So both disorder types break the symmetry somewhat;
      the inductor case simply does so more directly. The measured
      $1.38\times$ is the honest size of that effect for an 8-cell chain.
\item \textbf{Conclusion:} the disorder that breaks chiral symmetry is
      measurably more damaging --- isolating chiral symmetry as the
      \emph{reason} the edge state is robust. Practically: inductor quality
      matters more than capacitor quality.
\end{itemize}

% =========================================================================
% FIGURE 16
% =========================================================================
\clearpage
\section{Figure 16\quad Topological vs.\ trivial circuit at 10\% disorder}
\figblock{fig16_mc_topo_vs_trivial.png}
\theorybox{Swapping $C_1$ and $C_2$ moves the circuit into the
topologically \emph{trivial} phase, which has no protected gap state. The
bulk--boundary theorem then predicts \emph{no} robust resonance inside the
gap --- the gap should stay empty whether or not disorder is present.}

\shows
\begin{itemize}
\item Median impedance spectra under $10\%$ disorder for the topological
      circuit ($C_1<C_2$, blue) and the trivial circuit ($C_1>C_2$,
      green); dotted lines are the clean spectra, faint bands the spread.
\item At the gap centre (dashed line) the topological circuit shows a
      \emph{peak}; the trivial circuit shows a \emph{dip}.
\end{itemize}

\method
\begin{itemize}
\item For each phase, 400 disordered circuits were simulated across the
      full band and the median spectrum extracted.
\end{itemize}

\analysis
\begin{itemize}
\item \textbf{Quantitative:} the clean topological peak is
      $146$~k$\Omega$; the clean trivial circuit reaches only
      $48$~k$\Omega$ and not in the gap --- a factor of 3 smaller. Under
      $10\%$ disorder the topological median still bumps \emph{up} at the
      gap centre while the trivial median dips \emph{down}.
\item \textbf{Qualitative:} the trivial circuit's resonances are ordinary
      bulk modes with nothing protecting them; the gap is simply empty.
\item \textbf{Deviations:} the wide percentile bands are expected --- with
      $10\%$ disorder individual bulk resonances wander, so a
      \emph{point-by-point} band is broad; the \emph{median} still carries
      the message cleanly.
\item \textbf{Conclusion:} the robust gap resonance exists \emph{only} in
      the topological phase. This is the decisive control: the robustness
      documented in Figs.~9--15 is genuinely topological in origin, not a
      generic property of any LC ladder.
\end{itemize}

% =========================================================================
% SUMMARY
% =========================================================================
\clearpage
\section*{Summary}

The topological boundary resonance of the updated SSH circuit
($C_1=2.2$~nF, $C_2=3.3$~nF, $L=2200\ \mu$H, $1\ \Omega$ series resistance
throughout) sits at \textbf{45.74~kHz}, agreeing with analytic theory to
0.025\% --- the residual being a fully-explained finite-size effect. It
survives a single impurity (shift $\le2.6\%$), a clustered AB/BA impurity
pair (a few percent), and random $\pm10\%$ disorder on every capacitor: the
resonance stays centred on its pristine value, its peak height falls only
$3.6\%$, and the edge state keeps both its shape and its decay length.
Every deviation from ideal behaviour was traced to its cause --- finite
chain length, the $1/\sqrt{\;}$ convexity of the resonance formula, and the
expected second-order biases of a random ensemble. Inductor disorder, which
breaks chiral symmetry, is measurably more damaging than capacitor
disorder, and the trivial-phase control shows the robust gap resonance
exists \emph{only} in the topological phase. Every theory check passed.

% =========================================================================
% APPENDIX A -- SOURCES
% =========================================================================
\clearpage
\section*{Appendix A\quad Sources}
\small
The physics and circuit conventions reconstructed here are drawn from the
following sources.

\begin{enumerate}[leftmargin=1.6em,itemsep=3pt]
\item \textbf{Project brief:} \texttt{robustness\_manual.md} --- the
      step-by-step manual specifying the four-part study, the circuit
      topology and the component conventions reproduced in this report.
\item \textbf{Su, W.\,P., Schrieffer, J.\,R.\ \& Heeger, A.\,J.} (1979).
      Solitons in polyacetylene. \textit{Physical Review Letters}
      \textbf{42}, 1698 --- the original SSH model and its domain-wall
      soliton.
\item \textbf{Lee, C.\,H.\ \textit{et al.}} (2018). Topolectrical
      Circuits. \textit{Communications Physics} \textbf{1}, 39 --- the
      mapping of the SSH model onto an LC circuit and the grounded-Laplacian
      formalism.
\item \textbf{Anzurez, J.\ \textit{et al.}} (2025). Edge state robustness
      in the SSH model with impurities: topoelectric circuit simulations
      toward topological heterostructures. \textit{Physica B} \textbf{715},
      417609 --- the single-impurity and AB/BA domain-wall conventions
      (Parts 2--3).
\item \textbf{Asb\'oth, J.\,K., Oroszl\'any, L.\ \& P\'alyi, A.} (2016).
      \textit{A Short Course on Topological Insulators}. Lecture Notes in
      Physics \textbf{919}, Springer --- chiral symmetry, the bulk--boundary
      correspondence and the localization length.
\item \textbf{Anderson, P.\,W.} (1958). Absence of diffusion in certain
      random lattices. \textit{Physical Review} \textbf{109}, 1492 --- the
      theory of disorder-induced localization underpinning Part~4.
\item \textbf{Numerical tools:} the nodal-admittance method is standard
      linear circuit analysis; the solver here uses NumPy (Harris
      \textit{et al.}, \textit{Nature} \textbf{585}, 357, 2020) and
      Matplotlib (Hunter, \textit{Comput.\ Sci.\ Eng.}\ \textbf{9}, 90,
      2007). Figure colours follow the colour-blind-safe Okabe--Ito
      palette.
\end{enumerate}

% =========================================================================
% APPENDIX B -- RESULTS TABLE
% =========================================================================
\clearpage
\section*{Appendix B\quad Table of numerical results vs.\ theory}
\small
All values produced by this study's solver. ``Theory'' is the independent
analytic expectation; the final column is the agreement or the explanation
of any deviation.

\renewcommand{\arraystretch}{1.25}
\begin{longtable}{p{4.9cm}p{3.0cm}p{3.0cm}p{4.4cm}}
\toprule
\textbf{Quantity} & \textbf{Simulation} & \textbf{Theory} &
\textbf{Agreement / note}\\
\midrule
\endhead
\multicolumn{4}{l}{\textit{Part 1 --- pristine baseline}}\\
TBR frequency $f_R$        & $45.742$ kHz & $45.754$ kHz & 0.025\% (finite-size hybridisation)\\
TBR from Laplacian eigenvalue & $45.742$ kHz & $45.754$ kHz & matches simulation exactly\\
Peak impedance $|Z|_{\rm peak}$ & $146.3$ k$\Omega$ & --- & finite (set by $1\ \Omega$ losses)\\
Quality factor $Q$        & $892$ & --- & width $51$ Hz\\
Band gap                  & $41.8$--$51.2$ kHz & --- & 9.4 kHz wide\\
Decay length $\xi$        & $2.4656$ cells & $2.4663$ cells & 0.03\%\\
Per-cell voltage ratio    & $0.66658$ & $0.66667$ & 0.01\%\\
Exponential-fit quality   & $R^2=0.99999998$ & $R^2=1$ & near-perfect\\
Pristine IPR (N2--N16)    & $0.385$ & --- & vs.\ delocalized $0.067$\\
\midrule
\multicolumn{4}{l}{\textit{Part 2 --- single impurity}}\\
Largest shift, cell 2 $\delta=-0.3$ & $+2.57\%$ & --- & raises $f$ (smaller $C$)\\
Largest shift, cell 2 $\delta=+0.3$ & $-2.40\%$ & --- & lowers $f$ (larger $C$)\\
$\pm\delta$ asymmetry      & $\times1.07$ & $\times1$ (1st order) & $1/\sqrt{\;}$ convexity\\
Cell 2 / cell 7 leverage  & $\times51$ & --- & edge-mode amplitude\\
A-sublattice upstream change & $\le0.32\%$ & $0$ & defect cannot act upstream\\
\midrule
\multicolumn{4}{l}{\textit{Part 3 --- domain walls}}\\
AB pair shift span        & $3.80\%$ & --- & three capacitors changed\\
BA pair shift span        & $2.47\%$ & --- & BA disturbs TBR less\\
$f$-vs-$\gamma$ curvature (AB) & $+0.018$ & $0$ (1st order) & $1/\sqrt{\;}$ convexity\\
Deviation: sim vs.\ exact eigenvalue & $<0.001\%$ & --- & solver is exact\\
Deviation: sim vs.\ 1st-order (AB) & up to $0.91\%$ & --- & linearised-theory error\\
Domain-wall even-node IPR & AB $0.57$ / BA $0.38$ & --- & BA more uniform\\
\midrule
\multicolumn{4}{l}{\textit{Part 4 --- Monte Carlo ($M=8000$ per $\sigma$)}}\\
$\langle f\rangle$ at $\sigma=10\%$ & $45.898$ kHz & $45.742$ kHz & $+0.34\%$ bias (Jensen, $\propto\sigma^2$)\\
$f$ spread at $\sigma=1/5/10\%$ & $0.10/0.51/1.04$ kHz & $\propto\sigma$ & ratio $1.01,\ 1.03$\\
$\langle|Z|_{\rm peak}\rangle$ at $10\%$ & $141.0$ k$\Omega$ & $146.3$ k$\Omega$ & $-3.6\%$ bias (detuning, $\propto\sigma^2$)\\
$\langle$IPR$\rangle$ at $10\%$ & $0.355$ & $0.385$ & stays $5\times$ delocalized value\\
$\langle\xi\rangle$ at $10\%$ & $2.686$ cells & $2.466$ cells & $+8.9\%$ (bulk admixture)\\
Chiral test: $f$ spread at $10\%$ & cap $2.27\%$ / ind $3.13\%$ & ind $>$ cap & ind $\times1.38$ (chiral-breaking)\\
Trivial-phase clean peak  & $48$ k$\Omega$ & no gap mode & vs.\ topological $146$ k$\Omega$\\
\bottomrule
\end{longtable}

% =========================================================================
% APPENDIX C -- GLOSSARY
% =========================================================================
\clearpage
\section*{Appendix C\quad Glossary}
\small
Plain-language definitions of the technical terms used in this report.

\begin{description}[leftmargin=0pt,itemsep=3pt,style=unboxed]
\item[SSH model] The Su--Schrieffer--Heeger model: a one-dimensional chain
  with alternating ``strong'' and ``weak'' links. It is the simplest model
  with a topological edge state.
\item[Topolectrical circuit] An electronic circuit (inductors and
  capacitors) wired so that its equations are mathematically identical to a
  quantum model --- here the SSH model. It lets a quantum effect be
  measured on a bench.
\item[Unit cell] The repeating building block of the chain. This circuit
  has 8 unit cells, each containing 2 nodes.
\item[Node] A junction in the circuit where components meet. The chain has
  16 nodes; node voltages are what the simulation solves for.
\item[Intracell / intercell capacitor] $C_1$ links the two nodes
  \emph{within} a unit cell; $C_2$ links \emph{neighbouring} cells. Here
  $C_1<C_2$, which makes the circuit topological.
\item[Resonance / resonant frequency] A frequency at which a circuit
  responds very strongly. An LC loop resonates at $f=1/(2\pi\sqrt{LC})$.
\item[Topological boundary resonance (TBR)] The special resonance of this
  circuit that lives at the chain's edge and inside the band gap. It is the
  electrical version of the SSH topological edge state.
\item[Edge state] A resonance whose voltage is concentrated at the end
  (edge) of the chain and decays into the interior, instead of being spread
  evenly.
\item[Bulk / bulk mode / bulk band] ``Bulk'' means the interior of the
  chain. Bulk modes are ordinary resonances spread across the whole chain;
  their frequencies group into two continuous \emph{bulk bands}.
\item[Band gap] The frequency window between the two bulk bands, where no
  bulk mode exists. The TBR lives alone here.
\item[Topological vs.\ trivial phase] Two regimes of the circuit. The
  topological phase ($C_1<C_2$) has a protected edge state; the trivial
  phase ($C_1>C_2$) does not. Swapping the two capacitors switches between
  them.
\item[Chiral symmetry] A mathematical symmetry of the SSH model (the two
  sublattices are interchangeable in a specific sense). It is what
  \emph{protects} the edge state; breaking it makes the state fragile.
\item[Sublattice (A / B)] The chain's nodes split into two interleaved
  groups --- odd nodes (A) and even nodes (B). The edge state lives almost
  entirely on one of them.
\item[Bulk--boundary correspondence] The theorem stating that a
  topological bulk (here $C_1<C_2$) \emph{must} produce a protected edge
  state --- one mode in the gap.
\item[Domain wall / soliton] A point where the chain's strong/weak bond
  pattern reverses. It binds an extra localized state, called a soliton in
  the original SSH model.
\item[AB / BA pair] Two impurities placed on adjacent sites; ``AB'' and
  ``BA'' are the two possible orderings, which behave differently.
\item[Impedance / input impedance] The ratio of voltage to current ---
  how strongly a circuit ``resists'' an AC drive. It peaks sharply at a
  resonance.
\item[Nodal-admittance matrix] The grid of numbers ($16\times16$ here)
  encoding how every node connects to every other. Solving a linear system
  with this matrix gives all the node voltages --- the core of the
  simulation.
\item[Capacitance Laplacian (lossless)] A simplified version of the
  circuit matrix, ignoring resistive losses. Its \emph{eigenvalues} give
  the circuit's resonant frequencies directly --- used as an independent
  theory check.
\item[Eigenvalue / eigenvector] Special numbers (eigenvalues) and patterns
  (eigenvectors) of a matrix. Each eigenvalue of the capacitance Laplacian
  corresponds to one resonance; the eigenvector is that mode's voltage
  shape.
\item[Quality factor $Q$ / FWHM] $Q$ measures how sharp a resonance is
  (high $Q$ = narrow, clean peak). FWHM (full width at half maximum) is the
  width of the peak; $Q=f/\text{FWHM}$.
\item[Series resistance] The small unavoidable resistance ($1\ \Omega$
  here) inside every real capacitor and inductor. It makes resonance peaks
  finite rather than infinitely tall.
\item[Parabolic refinement] A trick to locate a peak more precisely than
  the frequency grid: fit a parabola through the three points around the
  highest sampled point and take the parabola's true maximum.
\item[Perturbation theory / first order] An approximate method: estimate
  how a small change shifts a result, keeping only the term proportional to
  the change (``first order''). Accurate for small changes, less so for
  large ones.
\item[Inverse participation ratio (IPR)] One number measuring how
  concentrated a voltage profile is. High IPR = localized on few nodes;
  low IPR = spread out. The fully-spread value here is $1/15\approx0.067$.
\item[Localization length $\xi$] The physical size of the edge state --- how
  many unit cells the voltage takes to fall by a factor $e$ ($\approx2.7$).
\item[Monte Carlo] A method that runs a simulation thousands of times with
  random inputs to build up a statistical distribution of outcomes.
\item[Gaussian disorder / $\sigma$] Each component is multiplied by a
  random number drawn from a bell curve (Gaussian) of width $\sigma$.
  $\sigma=10\%$ means a typical $\pm10\%$ manufacturing tolerance.
\item[Median / percentile / 5--95\% band] The median is the middle value
  of a sample. The 5--95\% band is the range containing the central 90\% of
  the random circuits --- a measure of spread.
\item[Bootstrap] A way to estimate the uncertainty of a statistic by
  re-sampling the data many times with replacement.
\item[Jensen's inequality / convexity] A curved (convex) function, such as
  $1/\sqrt{x}$, has the property that its \emph{average over a spread of
  inputs} is larger than its value at the average input. This produces the
  small upward bias of the mean resonance frequency under disorder.
\end{description}

\end{document}
